\section{Tautan GitHub}
\begin{itemize}
    \item Repository: \href{https://github.com/Vixrlie/KAI-Tubes-IF2224-2026}{https://github.com/Vixrlie/KAI-Tubes-IF2224-2026}
    \item Halaman release: \href{https://github.com/Vixrlie/KAI-Tubes-IF2224-2026/releases}{https://github.com/Vixrlie/KAI-Tubes-IF2224-2026/releases}
\end{itemize}

\section{Tabel Kontribusi}
\begin{center}
    \rowcolors{2}{gray!15}{white}
    \begin{tabularx}{\linewidth}{C{0.8cm} Y C{2.6cm} L C{1.6cm}}
        \toprule
        \textbf{No} & \textbf{Nama Lengkap} & \textbf{NIM} & \textbf{Deskripsi Pekerjaan} & \textbf{Kontribusi} \\
        \midrule
        1 & Jonathan Kris Wicaksono & 13524023 & AST: merancang node hierarchy AST, Syntax-Directed Translation untuk konversi parse tree ke AST, dan AST formatter untuk output Decorated AST. & 25\% \\
        2 & Vincent Rionarlie & 13524031 & Semantic Analyzer: implementasi visit functions untuk setiap node, type checking dan compatibility rules, inferensi tipe, anotasi node AST, dan control flow validation. & 25\% \\
        3 & Muhammad Aufar Rizqi Kusuma & 13524061 & Symbol Table: implementasi \texttt{tab}, \texttt{btab}, \texttt{atab}, manajemen scope push/pop, registrasi dan lookup identifier, serta inisialisasi predefined identifier. & 25\% \\
        4 & Bryan Pratama Putra Hendra & 13524067 & Laporan dan README: penulisan laporan lengkap dan README, penyusunan bagian teori, serta quality assurance. & 25\% \\
        \bottomrule
    \end{tabularx}
\end{center}

\section{Cara Menjalankan Program}
\begin{lstlisting}[basicstyle=\ttfamily\scriptsize,breaklines=true,frame=single,numbers=none]
make -B
make run INPUT=test/milestone-3/input-1.txt
make run-output INPUT=test/milestone-3/input-1.txt OUTPUT=output-m3.txt
\end{lstlisting}